% TOP_PROBLEM: {"id": 1958, "title": "Erdős Problem #149", "category_id": 3, "category": {"id": 3, "name": "graph_theory", "display_name": "Graph Theory", "description": "Problems involving graphs, networks, and their properties.", "slug": "graph-theory", "order_index": 3, "created_at": "2026-07-31T15:26:25.671Z"}, "status": "open"}
% FIRST_POSTED: null
% ENTRY_KIND: statement_only
% Imported from: unsolvedmath_additional_500.tex; UnsolvedMath ID 1958
% !TeX program = lualatex
% Compile twice: lualatex 1958.tex
% Self-contained: no external bibliography, images, or \input files.
% Numeric section labels and visible numbers are original UnsolvedMath IDs.
\documentclass[11pt,a4paper]{article}
\usepackage[margin=24mm,headheight=14pt]{geometry}
\usepackage{fontspec}
\setmainfont{Latin Modern Roman}
\setsansfont{Latin Modern Sans}
\setmonofont{Latin Modern Mono}
\usepackage{amsmath,amssymb,mathtools}
\usepackage{unicode-math}
\setmathfont{Latin Modern Math}
\usepackage{microtype}
\usepackage{enumitem,xurl,needspace}
\usepackage[dvipsnames]{xcolor}
\usepackage{hyperref}
\hypersetup{unicode=true,colorlinks=true,linkcolor=MidnightBlue,urlcolor=MidnightBlue,
 pdftitle={UnsolvedMath Problem 1958},pdfauthor={Source-annotated compilation},bookmarksnumbered=true}
\usepackage{bookmark,tocloft,titlesec,fancyhdr}
\setlength{\cftsecnumwidth}{4.2em}
\cftsetpnumwidth{2.5em}
\cftsetrmarg{4em}
\titleformat{\section}{\Large\bfseries\raggedright}{\thesection}{0.7em}{}
\pagestyle{fancy}\fancyhf{}
\fancyhead[L]{\small\sffamily Additional UnsolvedMath statements}
\fancyhead[R]{\small Original catalogue IDs}
\fancyfoot[C]{\thepage}
\setlength{\parindent}{0pt}
\setlength{\parskip}{5pt plus 1pt}
\setlength{\emergencystretch}{3em}
\setcounter{tocdepth}{1}\setcounter{secnumdepth}{1}
\setlist[itemize]{leftmargin=3em,itemsep=4pt,topsep=4pt,parsep=0pt}
\allowdisplaybreaks
\newcommand{\N}{\mathbb{N}}
\newcommand{\Z}{\mathbb{Z}}
\newcommand{\Q}{\mathbb{Q}}
\newcommand{\R}{\mathbb{R}}
\newcommand{\C}{\mathbb{C}}
\newcommand{\F}{\mathbb{F}}
\newcommand{\A}{\mathbb{A}}
\newcommand{\PP}{\mathbb{P}}
\newcommand{\E}{\mathbb{E}}
\newcommand{\eps}{\varepsilon}
\newcommand{\dd}{\,\mathrm d}
\newcommand{\sref}[2]{\hyperlink{source:#1:#2}{[S#2]}}
\newif\ifshowcatalogue
\showcataloguetrue
\newcommand{\cataloguescope}[1]{\ifshowcatalogue
 \paragraph{Statement recorded in the supplied source.}
 #1\fi}
\begin{document}
% UnsolvedMath ID 1958; source catalogue code EP-149

\setcounter{section}{1957}

\section{The Strong Edge-Colouring Bound}\label{1958}

{\small\sffamily UnsolvedMath ID 1958 · Graph Theory}

\subsection{Definitions and mathematical statement}

\paragraph{Shared notation.}Unless a section says otherwise, $\N=\{1,2,\ldots\}$,
$\Z$ is the ring of integers, $\Q,\R,\C$ are the rational, real, and complex
fields, $[N]=\{1,\ldots,\lfloor N\rfloor\}$, and $\log$ is the natural
logarithm. For real functions of a parameter tending to infinity,
$f\ll g$ and $f=O(g)$ mean $|f|\leq Cg$ eventually for some fixed $C$;
$f\gg g$ reverses this comparison; $f=o(g)$ means $f/g\to0$;
$f\sim g$ means $f/g\to1$; and $f\asymp g$ means both order bounds.
A subscript on $O$ or $\ll$ specifies allowed constant dependence.
The expression $n^{o(1)}$ in an upper bound means growth smaller than
$n^\epsilon$ for each fixed $\epsilon>0$. Local definitions govern any
exception, finite-field convention, or use of the same symbol for another
object. An asymptotic claim is not automatically an effective algorithm.



An induced matching is a set of vertex-disjoint edges with no further edge between their endpoints. A strong edge colouring partitions $E(G)$ into induced matchings. This is stricter than a proper edge colouring, which only separates edges sharing a vertex.

A graph has a vertex set $V$ and edges that are unordered pairs of distinct vertices. A subgraph need not be induced unless that word is stated. A proper vertex colouring gives adjacent vertices different colours; $\chi(G)$ is the least number of colours. \sref{1958}{1} \sref{1958}{2}

\paragraph{Statement recorded in the supplied source.}
Let $G$ be a graph with maximum degree $\Delta$. Is $G$ the union of at most $\tfrac{5}{4}\Delta^2$ sets of strongly independent edges (sets such that the induced subgraph is the union of vertex-disjoint edges)? \sref{1958}{1}

\subsection{Short English statement}

Can every graph's edges be split into induced matchings using at most five quarters of the square of its maximum degree?

\subsection{Sources}

{\small

\begin{itemize}

\item[\hypertarget{source:1958:1}{S1}] ULAM.AI, \emph{UnsolvedMath}, supplied \texttt{problems.json}, record 1958 (EP-149), statement and background. The embedded literature-review date is 2026-08-17; that is the archive’s review date, not a fresh certification. Dataset landing page: \url{https://huggingface.co/datasets/ulamai/UnsolvedMath}.

\item[\hypertarget{source:1958:2}{S2}] T. F. Bloom, \emph{Erdős Problems}, Problem 149, \url{https://www.erdosproblems.com/149}. Reference imported from the supplied record; the live problem page was not independently checked for this compilation.

\item[\hypertarget{source:1958:3}{S3}] Eoin Hurley, Remi de Joannis de Verclos, and Ross J. Kang, An improved procedure for colouring graphs of bounded local density, Advances in Combinatorics (2022), Paper 7, \nolinkurl{doi:10.19086/aic.2022.7.} \url{https://arxiv.org/abs/2007.07874}. Bibliographic information imported from the supplied record.

\item[\hypertarget{source:1958:4}{S4}] Wouter Cames van Batenburg, Ross J. Kang, and Francois Pirot, Strong cliques and forbidden cycles, Indagationes Mathematicae 31 (2020), 64-82, \nolinkurl{doi:10.1016/j.indag.2019.09.003.} \url{https://arxiv.org/abs/1903.06087}. Bibliographic information imported from the supplied record.

\item[\hypertarget{source:1958:5}{S5}] Bonamy, Marthe and Perrett, Thomas and Postle, Luke, Colouring graphs with sparse neighbourhoods: bounds and applications. J. Combin. Theory Ser. B (2022), 278-317. Bibliographic information imported from the supplied record.

\item[\hypertarget{source:1958:6}{S6}] Bruhn, Henning and Joos, Felix, A stronger bound for the strong chromatic index. Combin. Probab. Comput. (2018), 21-43. Bibliographic information imported from the supplied record.

\end{itemize}

}

\label{end:1958}
\end{document}
